\begin{definition}[Parent ground space under the canonical $\ell^2$ identification]
    \label{def:parent_ground_space_es}
    \lean{MPSTensor.parentHamiltonianGroundSpaceES}
    \leanok
    \uses{def:parent_hamiltonian, def:nsite_space_parent}
    Denote by $\mathcal G_{N,L}^{\mathrm{ES}}(A)$ the parent ground space after
    identifying the coefficient space with the Hilbert space
    $\ell^{2}(\{0,\ldots,d-1\}^N)$.
\end{definition}

\begin{definition}[Parent Hamiltonian under the canonical $\ell^2$ identification]
    \label{def:parent_hamiltonian_es}
    \lean{MPSTensor.parentHamiltonianES}
    \leanok
    \uses{def:parent_hamiltonian, def:nsite_space_parent}
    Denote by $H_{N,L}^{\mathrm{ES}}(A)$ the conjugate of $H_{N}(A,L)$ by the
    canonical identification between the coefficient space and
    $\ell^{2}(\{0,\ldots,d-1\}^N)$.
\end{definition}

\begin{theorem}[Kernel under the canonical $\ell^2$ identification]
    \label{thm:parent_ground_space_es_kernel}
    \lean{MPSTensor.parentHamiltonianGroundSpaceES_eq_ker_parentHamiltonianES}
    \leanok
    Under the canonical $\ell^2$ identification of the $N$-site state space, the
    parent ground space is exactly the kernel of the corresponding
    parent Hamiltonian.
\end{theorem}

\begin{proof}
    \leanok
    This is an immediate unraveling of the definitions: both constructions are
    obtained from the same parent Hamiltonian by conjugating with the canonical
    identification between the site space and the $\ell^2$ space, so
    the ground space is precisely the kernel after this identification.
\end{proof}

\begin{lemma}[Membership criterion under the canonical $\ell^2$ identification]
    \label{lem:mem_parent_ground_space_es}
    \lean{MPSTensor.mem_parentHamiltonianGroundSpaceES_iff}
    \leanok
    \uses{thm:parent_ground_space_es_kernel}
    A vector belongs to $\mathcal G_{N,L}^{\mathrm{ES}}(A)$ exactly when the
    parent Hamiltonian annihilates it.
\end{lemma}

\begin{remark}[Local projectors under the canonical $\ell^2$ identification]
    \label{rem:euclidean_local_projector_ingredients}
    \lean{MPSTensor.parentInteractionES}
    \lean{MPSTensor.cyclicRestrictES}
    \lean{MPSTensor.localTermES}
    \lean{MPSTensor.localTermESSummand}
    \leanok
    \uses{def:ground_space_es, rem:cyclic_window_operators}
    Let $P_L^{\mathrm{ES}}(A)$ denote the orthogonal projector onto
    $(G_L^{\mathrm{ES}}(A))^\perp$ on the canonical $\ell^2$ $L$-site Hilbert
    space. For a cyclic window starting at $i$ and an outside configuration
    $\tau$, the map $R_{i,\tau}$ restricts an $N$-site vector to that
    length-$L$ cyclic window with the outside configuration fixed by $\tau$.
    The local term
    $h_{i,\mathrm{ES}}$ is the corresponding
    $N$-site local projector, and $R_{i,\tau}^{\dagger} P_L^{\mathrm{ES}}(A) R_{i,\tau}$ is
    the basic positive summand in its cyclic-averaging formula.
\end{remark}

\begin{lemma}[Positivity of the Hilbert-space local summands]
    \label{lem:euclidean_local_projector_positive}
    \lean{MPSTensor.parentInteractionES_isPositive}
    \lean{MPSTensor.localTermESSummand_isPositive}
    \lean{MPSTensor.localTermESSummand_sum_isPositive}
    \leanok
    \uses{rem:euclidean_local_projector_ingredients}
    The Hilbert-space local projector $P_L^{\mathrm{ES}}(A)$ is positive.
    Consequently, every summand
    $R_{i,\tau}^{\dagger} P_L^{\mathrm{ES}}(A) R_{i,\tau}$ is positive, and so is
    the finite sum over $\tau$.
\end{lemma}

\begin{proof}
    \leanok
    Orthogonal projectors are positive. Conjugation by $R_{i,\tau}$ preserves
    positivity, and finite sums of positive operators remain positive.
\end{proof}

\begin{lemma}[Cyclic averaging for the Hilbert-space local terms]
    \label{lem:local_term_es_average}
    \lean{MPSTensor.parentHamiltonianES_eq_sum_localTermES}
    \lean{MPSTensor.localTermES_eq_average_localTermESSummand}
    \leanok
    \uses{rem:euclidean_local_projector_ingredients, lem:euclidean_local_projector_positive}
    The Hilbert-space parent Hamiltonian is the sum of the Hilbert-space
    local terms $H_{N,\mathrm{ES}} = \sum_i h_{i,\mathrm{ES}}$, and each local
    term admits the exact cyclic-averaging formula
    \begin{equation}\label{eq:ph_cyclic_averaging}
        h_{i,\mathrm{ES}}
        = d^{-L} \sum_{\tau} R_{i,\tau}^{\dagger} P_L^{\mathrm{ES}}(A) R_{i,\tau}.
    \end{equation}
\end{lemma}

\begin{proof}
    \leanok
    Unfold the Hilbert-space local term and compare it
    configuration-by-configuration with the cyclic restrictions. The $d^{-L}$
    factor compensates for the $d^L$ choices of the window coordinates in the
    ambient configuration parameter $\tau$.
\end{proof}

\begin{lemma}[Positivity of the Hilbert-space Hamiltonian]
    \label{lem:transported_es_positive}
    \lean{MPSTensor.localTermES_isPositive}
    \lean{MPSTensor.parentHamiltonianES_isPositive}
    \leanok
    \uses{lem:local_term_es_average}
    Each Hilbert-space local term $h_{i,\mathrm{ES}}$ is positive, hence the full
    Hilbert-space parent Hamiltonian $H_{N,\mathrm{ES}}$ is positive.
\end{lemma}

\begin{proof}
    \leanok
    Apply the averaging formula~(\ref{eq:ph_cyclic_averaging}):
    $h_{i,\mathrm{ES}}$ is a non-negative scalar multiple of a finite sum of
    positive conjugates of $P_L^{\mathrm{ES}}(A)$, hence is itself positive.
    Summing over $i$ yields positivity of $H_{N,\mathrm{ES}}$.
\end{proof}

\begin{lemma}[Hilbert-space local terms are symmetric projections]
    \label{lem:transported_es_local_projections}
    \lean{MPSTensor.parentInteractionES_isSymmetricProjection}
    \lean{MPSTensor.localTermES_isSymmetricProjection}
    \leanok
    \uses{rem:euclidean_local_projector_ingredients, lem:transported_es_positive}
    The Hilbert-space local projector $P_L^{\mathrm{ES}}(A)$ is a symmetric
    projection, and every Hilbert-space local term $h_{i,\mathrm{ES}}$ on the
    $N$-site chain is a symmetric projection.
\end{lemma}

\begin{proof}
    \leanok
    \uses{lem:transported_es_positive}
    The local operator $P_L^{\mathrm{ES}}(A)$ is an orthogonal projection onto
    $(G_L^{\mathrm{ES}}(A))^\perp$. For $L\le N$, restricting
    $h_{i,\mathrm{ES}}v$ to a cyclic window gives
    $P_L^{\mathrm{ES}}(A)$ applied to the corresponding restriction of $v$; applying
    $h_{i,\mathrm{ES}}$ a second time therefore applies
    $(P_L^{\mathrm{ES}}(A))^2=P_L^{\mathrm{ES}}(A)$ on that window. The positivity
    result above supplies symmetry, and the case $L>N$ is the zero projection by
    definition.
\end{proof}

\begin{lemma}[Kernel of the Hilbert-space parent interaction]
    \label{lem:parent_interaction_es_kernel}
    \lean{MPSTensor.parentInteractionES_apply_eq_zero_iff}
    \leanok
    \uses{def:ground_space_es, rem:euclidean_local_projector_ingredients}
    The Hilbert-space local interaction $P_L^{\mathrm{ES}}(A)$ annihilates exactly
    the Hilbert-space local ground space:
    \[
        P_L^{\mathrm{ES}}(A)v=0 \quad\Longleftrightarrow\quad
        v\in G_L^{\mathrm{ES}}(A).
    \]
\end{lemma}

\begin{proof}
    \leanok
    Since $P_L^{\mathrm{ES}}(A)$ is the orthogonal projector onto
    $(G_L^{\mathrm{ES}}(A))^\perp$, its kernel is the original subspace
    $G_L^{\mathrm{ES}}(A)$.
\end{proof}

\begin{lemma}[Local-term kernel via cyclic restrictions]
    \label{lem:local_term_es_kernel_restrictions}
    \lean{MPSTensor.localTermES_eq_zero_iff_forall_cyclicRestrictES_mem_groundSpaceES}
    \lean{MPSTensor.cyclicRestrictES_mem_groundSpaceES_of_localTermES_eq_zero}
    \leanok
    \uses{lem:parent_interaction_es_kernel, rem:euclidean_local_projector_ingredients}
    For $L\le N$, a Hilbert-space local term $h_{i,\mathrm{ES}}$ annihilates a
    vector $v$ iff every restriction of $v$ to the cyclic length-$L$ window
    starting at $i$, with the outside configuration fixed, lies in
    $G_L^{\mathrm{ES}}(A)$.
\end{lemma}

\begin{proof}
    \leanok
    Restricting $h_{i,\mathrm{ES}}v$ to the same cyclic window gives
    $P_L^{\mathrm{ES}}(A)$ applied to the corresponding restriction of $v$.
    Lemma~\ref{lem:parent_interaction_es_kernel} identifies the vanishing of this
    local projector with membership in $G_L^{\mathrm{ES}}(A)$.
\end{proof}

\begin{theorem}[Zero energy gives the cyclic local constraints]
    \label{thm:parent_hamiltonian_kernel_le_chain_ground_space}
    \lean{MPSTensor.ker_parentHamiltonian_le_chainGroundSpace}
    \leanok
    \uses{def:parent_hamiltonian, def:parent_hamiltonian_es,
        def:chain_ground_space, lem:transported_es_local_projections,
        lem:local_term_es_kernel_restrictions, rem:projection_geometry_rowsum_identities}
    Assume $N>0$ and $L\le N$. If a vector $\psi$ has zero energy for the
    finite periodic parent Hamiltonian,
    \[
        H_N(A,L)\psi=0,
    \]
    then every cyclic length-$L$ restriction of $\psi$ lies in the local MPS
    ground space. Equivalently,
    \[
        \ker H_N(A,L) \subseteq G_L^{(N)}(A).
    \]
\end{theorem}

\begin{proof}
    \leanok
    Let $e_N$ be the canonical $\ell^2$ identification of the $N$-site space.
    Since $H_N(A,L)$ is the $e_N$-conjugate of $H_{N,L}^{\mathrm{ES}}(A)$, the
    hypothesis gives
    \[
        H_{N,L}^{\mathrm{ES}}(A)(e_N^{-1}\psi)
        =
        e_N^{-1}\bigl(H_N(A,L)\psi\bigr)
        =
        0.
    \]
    Put $P_i=h_{i,\mathrm{ES}}$ and $H=\sum_iP_i$. If $Hv=0$, then
    \[
        0
        =
        \re\langle Hv,v\rangle
        =
        \sum_i\re\langle P_i v,v\rangle
        =
        \sum_i\|P_i v\|^2 .
    \]
    Each summand is nonnegative; hence $\|P_i v\|^2=0$ and therefore
    $P_i v=0$ for every $i$. Applying this to
    $H_{N,L}^{\mathrm{ES}}(A)=\sum_i h_{i,\mathrm{ES}}$ gives the vanishing of
    every Hilbert-space local term. Lemma~\ref{lem:local_term_es_kernel_restrictions}
    identifies these equations with the conditions: for all $i$ and $\tau$,
    \[
        R_{i,\tau}\psi\in G_L(A),
    \]
    which are exactly the cyclic-window defining equations of
    $G_L^{(N)}(A)$.
\end{proof}

\begin{lemma}[Cyclic constraints give zero energy]
    \label{lem:chain_ground_space_le_parent_hamiltonian_kernel}
    \lean{MPSTensor.chainGroundSpace_le_ker_parentHamiltonian}
    \leanok
    \uses{def:parent_hamiltonian, def:chain_ground_space,
        thm:chain_ground_space_implies_frustration_free}
    Assume $N>0$ and $L\le N$. Then
    \[
        G_L^{(N)}(A)\subseteq \ker H_N(A,L).
    \]
\end{lemma}

\begin{proof}
    \leanok
    If \(\psi\in G_L^{(N)}(A)\), then every local parent interaction
    annihilates \(\psi\), for every site \(i\):
    \[
        h_i(A,L)\psi=0.
    \]
    Summing over all sites gives
    \[
        H_N(A,L)\psi=\sum_i h_i(A,L)\psi=0.
    \]
\end{proof}

\begin{definition}[Hilbert-space cyclic constraint space]
    \label{def:chain_ground_space_es}
    \lean{MPSTensor.chainGroundSpaceES}
    \leanok
    \uses{def:chain_ground_space}
    The Hilbert-space cyclic constraint space is the transport of
    \(G_L^{(N)}(A)\) by the canonical \(\ell^2\)-identification:
    \[
        G_{L,\mathrm{ES}}^{(N)}(A)=e_N^{-1}G_L^{(N)}(A).
    \]
\end{definition}

\begin{theorem}[The finite-volume kernel is the cyclic constraint space]
    \label{thm:parent_hamiltonian_kernel_eq_chain_ground_space}
    \lean{MPSTensor.ker_parentHamiltonian_eq_chainGroundSpace}
    \lean{MPSTensor.parentHamiltonianGroundSpaceES_eq_chainGroundSpaceES}
    \leanok
    \uses{def:parent_hamiltonian, def:parent_hamiltonian_es,
        def:chain_ground_space, thm:parent_hamiltonian_kernel_le_chain_ground_space,
        lem:chain_ground_space_le_parent_hamiltonian_kernel,
        def:chain_ground_space_es}
    Assume $N>0$ and $L\le N$. Then
    \[
        \ker H_N(A,L)=G_L^{(N)}(A).
    \]
    Equivalently, after the canonical Hilbert-space identification,
    \[
        \mathcal K_N^{\mathrm{ES}}(A,L)
        =
        e_N^{-1}G_L^{(N)}(A),
    \]
    where \(\mathcal K_N^{\mathrm{ES}}(A,L)\) is the zero-energy space of
    \(H_{N,L}^{\mathrm{ES}}(A)\).
\end{theorem}

\begin{proof}
    \leanok
    The inclusion
    \[
        \ker H_N(A,L)\subseteq G_L^{(N)}(A)
    \]
    is Theorem~\ref{thm:parent_hamiltonian_kernel_le_chain_ground_space}. For
    the converse, a vector in \(G_L^{(N)}(A)\) satisfies all local equations
    \(h_i(A,L)\psi=0\). Therefore
    \[
        H_N(A,L)\psi=\sum_i h_i(A,L)\psi=0.
    \]
    The Hilbert-space statement is the same equality transported by \(e_N^{-1}\).
\end{proof}

\begin{theorem}[Intersection property for two adjacent local kernels]
    \label{thm:adjacent_local_kernels_ground_space_es}
    \lean{MPSTensor.mem_groundSpaceES_succ_of_adjacent_localTermES_eq_zero}
    \lean{MPSTensor.adjacent_localTermES_eq_zero_of_mem_groundSpaceES_succ}
    \lean{MPSTensor.adjacent_localTermES_eq_zero_iff_mem_groundSpaceES_succ}
    \leanok
    \uses{thm:ground_space_intersection, lem:local_term_es_kernel_restrictions,
        lem:mem_ground_space_es}
    If $A$ is injective and $L\ge2$, then on an $(L+1)$-site chain the kernels of
    the two adjacent Hilbert-space $L$-site local terms have common kernel equal
    to the Hilbert-space $(L+1)$-site MPS ground space:
    \[
        h_{0,\mathrm{ES}}v=0\ \text{ and }\ h_{1,\mathrm{ES}}v=0
        \quad\Longleftrightarrow\quad
        v\in G_{L+1}^{\mathrm{ES}}(A).
    \]
\end{theorem}

\begin{proof}
    \leanok
    \uses{thm:ground_space_intersection, lem:local_term_es_kernel_restrictions,
        lem:mem_ground_space_es}
    Lemma~\ref{lem:local_term_es_kernel_restrictions} translates the two local
    kernel assumptions into the left and right $L$-site ground conditions. The
    open-chain intersection property, Theorem~\ref{thm:ground_space_intersection},
    then grows back the shared $(L-1)$-site overlap to the $(L+1)$-site ground space.
    The converse uses the forward restriction lemmas for ground-space vectors and
    the same kernel--restriction characterization.
\end{proof}

\begin{lemma}[Non-overlap positivity for Hilbert-space local terms]
    \label{lem:transported_es_nonoverlap_positive}
    \lean{MPSTensor.CyclicWindowsDisjoint}
    \lean{MPSTensor.CyclicWindowsDisjoint.symm}
    \lean{MPSTensor.CyclicWindowsDisjoint.of_not_cyclicWindowsOverlap}
    \lean{MPSTensor.localTerm_commute_of_cyclic_windows_disjoint}
    \lean{MPSTensor.localTermES_commute_of_cyclic_windows_disjoint}
    \lean{MPSTensor.localTermES_re_inner_nonneg_of_cyclic_windows_disjoint}
    \leanok
    \uses{lem:transported_es_local_projections, rem:projection_geometry_rowsum_identities}
    Let $L\le N$ and let $h_{i,\mathrm{ES}},h_{j,\mathrm{ES}}$ be Hilbert-space local
    terms whose cyclic $L$-windows are site-disjoint: no site of the periodic chain
    has cyclic offset $<L$ from both $i$ and $j$. Then the terms commute pointwise:
    for every vector $v$,
    \[
        h_{i,\mathrm{ES}}(h_{j,\mathrm{ES}}v)
        = h_{j,\mathrm{ES}}(h_{i,\mathrm{ES}}v),
    \]
    The same commutation identity holds for the corresponding local terms before
    passing to the Hilbert-space model. Their ordered cross term is non-negative:
    \[
        0\le \re\langle h_{i,\mathrm{ES}}v,
        h_{j,\mathrm{ES}}v\rangle.
    \]
\end{lemma}

\begin{proof}
    \leanok
    \uses{lem:transported_es_local_projections, rem:projection_geometry_rowsum_identities}
    In coordinates, $h_{i,\mathrm{ES}}$ applies $P_L^{\mathrm{ES}}(A)$ only to the
    $i$-window variables and leaves the complementary variables fixed. Site
    disjointness gives two replacement identities: replacing the $i$-window does
    not change the extracted $j$-window, and the $i$- and $j$-window replacements
    commute. Therefore the two embedded local operators commute. The preceding
    projection-geometry identity for commuting symmetric projections gives the
    non-negative ordered cross term.
\end{proof}

\begin{lemma}[Quadratic form to norm bound]\label{lem:quadratic_form_to_norm_bound}
    \lean{FrustrationFree.spectralGap_of_martingale}
    \leanok
    Let $H$ be a positive semidefinite self-adjoint operator on a
    finite-dimensional complex Hilbert space, and suppose that, for all~$v$,
    \[
        \gamma\,\langle v, H v\rangle \le \langle H v, H v\rangle
    \]
    i.e.\ $H^2 \ge \gamma H$ as a quadratic form, for some $\gamma > 0$.
    Then for every $v$ in the orthogonal complement of $\ker H$,
    \[
        \gamma\,\|v\| \le \|H v\|.
    \]
    In particular, every positive eigenvalue of $H$ is at least $\gamma$.
    The positive-semidefiniteness hypothesis is essential: without it,
    $H = -\mathrm{Id}$ with $\gamma = 2$ satisfies the quadratic-form
    inequality vacuously but violates the norm bound.
\end{lemma}

\begin{proof}
    \leanok
    By the finite-dimensional spectral theorem, $H$ diagonalises in an
    orthonormal eigenbasis with real eigenvalues $\lambda_k$. For a unit
    eigenvector $\psi_k$ with eigenvalue $\lambda_k$ the hypothesis
    gives $\gamma\,\lambda_k \le \lambda_k^2$, whence $\lambda_k \le 0$
    or $\lambda_k \ge \gamma$. Since $H$ is positive semidefinite
    (its quadratic form is $\ge 0$), the negative case is excluded and
    every nonzero eigenvalue satisfies $\lambda_k \ge \gamma$. Writing
    $v = \sum_k c_k \psi_k$ and projecting onto $(\ker H)^{\perp}$
    keeps only terms with $\lambda_k \ge \gamma$, so
    $\|H v\|^2 = \sum_k |c_k|^2 \lambda_k^2
        \ge \gamma^2 \sum_k |c_k|^2 = \gamma^2 \|v\|^2$.
\end{proof}

\subsection{Martingale condition for a finite family of projections}\label{subsec:spectral_gap_rowsum}

\begin{lemma}[Parent-Hamiltonian quadratic-form reduction]
    \label{lem:parent_hamiltonian_es_quadratic_reduction}
    \lean{MPSTensor.parentHamiltonianES_norm_bound_of_quadratic_form}
    \lean{MPSTensor.parentHamiltonianES_gap_bound_of_quadratic_form}
    \leanok
    \uses{lem:transported_es_positive, thm:parent_ground_space_es_kernel,
        lem:quadratic_form_to_norm_bound}
    Let $H_{N,L}^{\mathrm{ES}}(A)$ be the Hilbert-space representative of the
    parent Hamiltonian. If a constant $\gamma > 0$ satisfies the uniform
    quadratic-form estimate
    \begin{equation}\label{eq:ph_uniform_quadratic_form}
        \gamma\,\re\langle H_{N,L}^{\mathrm{ES}}(A)v, v\rangle
        \le
        \re\langle H_{N,L}^{\mathrm{ES}}(A)v,
        H_{N,L}^{\mathrm{ES}}(A)v\rangle
    \end{equation}
    for every $N \ge 2L$ and every vector $v$, then
    \[
        \gamma\,\|v\| \le \|H_{N,L}^{\mathrm{ES}}(A)v\|
    \]
    for every $v \in (\mathcal G_{N,L}^{\mathrm{ES}}(A))^\perp$. In particular,
    with $L>1$ and $\gamma = 1/(4L)$, the explicit gap-bound conclusion follows
    once this same quadratic-form estimate~(\ref{eq:ph_uniform_quadratic_form}) is supplied.
\end{lemma}

\begin{proof}
    \leanok
    The Hilbert-space Hamiltonian is positive by
    Lemma~\ref{lem:transported_es_positive}, and its Hilbert-space ground
    space is its kernel by Theorem~\ref{thm:parent_ground_space_es_kernel}.
    Therefore the abstract spectral-theorem conversion of
    Lemma~\ref{lem:quadratic_form_to_norm_bound} applies directly. The
    explicit constant is positive because $L>1$ implies $4L>0$.
\end{proof}

\begin{remark}[Projection-geometry identities for row-sum arguments]
    \label{rem:projection_geometry_rowsum_identities}
    \lean{LinearMap.IsSymmetricProjection.re_inner_apply_apply_self}
    \lean{LinearMap.IsSymmetricProjection.re_inner_apply_apply_nonneg_of_commute}
    \lean{LinearMap.IsSymmetricProjection.re_inner_apply_apply_ge_neg_of_norm_apply_le}
    \lean{LinearMap.IsSymmetricProjection.re_inner_anticommutator_ge_neg_of_norm_apply_le}
    \lean{ProjectionGeometry.re_inner_sum_apply_left}
    \lean{ProjectionGeometry.apply_eq_zero_of_sum_apply_eq_zero}
    \lean{ProjectionGeometry.re_inner_sum_apply_apply}
    \lean{ProjectionGeometry.sum_sum_eq_diag_add_offdiag}
    \lean{ProjectionGeometry.crossTerm_sum_bound_of_ordered_rowSum}
    \lean{ProjectionGeometry.crossTerm_sum_bound_of_anticommutator_coeffSum}
    \lean{ProjectionGeometry.crossTerm_sum_bound_of_anticommutator_rowCol}
    \leanok
    The projection-geometry reduction uses the following elementary Hilbert-space
    identities: the diagonal identity
    $\re\langle Pv,Pv\rangle=\re\langle Pv,v\rangle$
    for symmetric projections; nonnegativity of this quadratic form; nonnegativity
    of $\re\langle Pv,Qv\rangle$ for commuting symmetric projections
    $P,Q$; the Cauchy--Schwarz conversion from
    $\|P Q v\|\le c\|Pv\|$ to
    $\re\langle Pv,Qv\rangle\ge -c\re\langle Pv,v\rangle$;
    the symmetric variant which, for $c\ge0$, converts the operator-product bound
    $\|P Q v\|\le c\|Qv\|$ to the anticommutator estimate
    $\re\langle Pv,Qv\rangle+\re\langle Qv,Pv\rangle
    \ge -c\bigl(\re\langle Pv,v\rangle
    +\re\langle Qv,v\rangle\bigr)$;
    finite-sum expansions of
    $\re\langle (\sum_i P_i)v,v\rangle$ and
    $\re\langle (\sum_i P_i)v,(\sum_i P_i)v\rangle$; the
    conclusion that a vector annihilated by a finite sum of symmetric projections
    is annihilated by each projection; the diagonal/off-diagonal split of the
    ordered double sum; the aggregate off-diagonal estimate obtained from row
    sums $\sum_{j\ne i}c_{ij}\le1$; and the corresponding aggregate estimate for
    the anticommutator martingale condition.
\end{remark}

\begin{lemma}[Ordered projection row-sum quadratic-form reduction]
    \label{lem:ordered_projection_rowsum_reduction}
    \lean{ProjectionGeometry.quadraticForm_sum_projections_of_ordered_rowSum}
    \leanok
    \uses{rem:projection_geometry_rowsum_identities}
    Let $P_i$ be a finite family of symmetric projections and let
    $H=\sum_i P_i$. If the ordered off-diagonal forms obey
    \[
        \re\langle P_i v, P_j v\rangle
        \ge -(1-\gamma)c_{ij}\re\langle P_i v,v\rangle
        \qquad (i\ne j),
    \]
    with row sums $\sum_{j\ne i} c_{ij}\le 1$ and $\gamma\le 1$, then
    $H^2\ge \gamma H$ as a quadratic form.
\end{lemma}

\begin{proof}
    \leanok
    Expand $\re\langle H v,Hv\rangle$ as an ordered double sum.
    Symmetric idempotence identifies each diagonal term with
    $\re\langle P_i v,v\rangle$, while the row-summable ordered
    cross-term bounds control the off-diagonal contribution by
    $-(1-\gamma)\sum_i\re\langle P_i v,v\rangle$.
\end{proof}

\begin{lemma}[Projection anticommutator row-sum reduction]
    \label{lem:anticommutator_projection_rowsum_reduction}
    \lean{ProjectionGeometry.quadraticForm_sum_projections_of_anticommutator_rowCol}
    \leanok
    \uses{rem:projection_geometry_rowsum_identities}
    Let $P_i$ be a finite family of symmetric projections and let
    $H=\sum_iP_i$. Suppose the anticommutator forms obey
    \[
        \re\langle P_i v,P_jv\rangle
        +\re\langle P_j v,P_iv\rangle
        \ge
        -(1-\gamma)c_{ij}
        \bigl(\re\langle P_i v,v\rangle
        +\re\langle P_j v,v\rangle\bigr)
        \qquad (i\ne j),
    \]
    with row sums $\sum_{j\ne i}c_{ij}\le1$, column sums
    $\sum_{i\ne j}c_{ij}\le1$, and $\gamma\le1$. Then
    $H^2\ge\gamma H$ as a quadratic form.
\end{lemma}

\begin{proof}
    \leanok
    Expand $\re\langle Hv,Hv\rangle$ as an ordered double sum. The
    anticommutator bound controls the sum of each ordered term with its reversed
    term. The row and column estimates bound the coefficient-weighted diagonal
    contribution by twice $\sum_i\re\langle P_i v,v\rangle$; after
    dividing by two, the same quadratic-form estimate follows.
\end{proof}

\begin{lemma}[Finite-overlap projection anticommutator reduction]
    \label{lem:finite_overlap_projection_anticommutator_reduction}
    \lean{ProjectionGeometry.quadraticForm_sum_projections_of_finite_overlap_anticommutator}
    \leanok
    \uses{lem:anticommutator_projection_rowsum_reduction}
    Let $P_i$ be a finite family of symmetric projections and let
    $\gamma\le1$. Suppose an interaction predicate $i\sim j$ has at most $m$
    off-diagonal neighbours in each row and in each column, where $m>0$. If
    noninteracting pairs have non-negative anticommutator form and interacting
    pairs satisfy
    \[
        \re\langle P_i v,P_jv\rangle
        +\re\langle P_j v,P_iv\rangle
        \ge
        -(1-\gamma)\frac{1}{m}
        \bigl(\re\langle P_i v,v\rangle
        +\re\langle P_j v,v\rangle\bigr),
    \]
    then $H=\sum_iP_i$ satisfies $H^2\ge\gamma H$ as a quadratic form.
\end{lemma}

\begin{proof}
    \leanok
    Choose $c_{ij}=1/m$ on interacting pairs and $c_{ij}=0$ otherwise. The row
    and column cardinality bounds give the corresponding coefficient sums, and
    the noninteracting anticommutator non-negativity supplies the zero-coefficient
    estimates. Apply Lemma~\ref{lem:anticommutator_projection_rowsum_reduction}.
\end{proof}

\begin{lemma}[Finite-overlap projection row-sum reduction]
    \label{lem:finite_overlap_projection_rowsum_reduction}
    \lean{ProjectionGeometry.quadraticForm_sum_projections_of_finite_overlap}
    \leanok
    \uses{lem:ordered_projection_rowsum_reduction}
    Let $P_i$ be a finite family of symmetric projections and let $\gamma\le1$.
    Suppose an interaction predicate $i\sim j$ has at most $m$ off-diagonal
    neighbours in each row for some positive integer $m$. If noninteracting
    ordered cross terms are non-negative and, on each interacting pair,
    \[
        \re\langle P_i v,P_jv\rangle
        \ge -(1-\gamma)\frac{1}{m}\re\langle P_i v,v\rangle,
    \]
    then $H=\sum_i P_i$ satisfies $H^2\ge \gamma H$ as a quadratic form.
\end{lemma}

\begin{proof}
    \leanok
    Choose $c_{ij}=1/m$ on interacting pairs and $c_{ij}=0$ otherwise. The row
    cardinality bound gives $\sum_{j\ne i}c_{ij}\le1$; noninteracting
    nonnegativity supplies the zero-coefficient cross-term estimates. Apply
    Lemma~\ref{lem:ordered_projection_rowsum_reduction}.
\end{proof}

\begin{lemma}[Finite-overlap projection norm-compression reduction]
    \label{lem:finite_overlap_projection_norm_reduction}
    \lean{ProjectionGeometry.quadraticForm_sum_projections_of_finite_overlap_norm_bound}
    \lean{ProjectionGeometry.quadraticForm_sum_projections_of_finite_overlap_norm_bound_of_le}
    \lean{ProjectionGeometry.quadraticForm_sum_projections_of_finite_overlap_norm_bound_of_lt}
    \leanok
    \uses{lem:finite_overlap_projection_rowsum_reduction,
        rem:projection_geometry_rowsum_identities}
    Let $P_i$ be a finite family of symmetric projections and let $\gamma\le1$.
    Suppose an interaction predicate has at most $m>0$ off-diagonal neighbours
    in each row, noninteracting ordered cross terms are non-negative, and
    interacting pairs satisfy the norm-compression estimate
    \begin{equation}\label{eq:ph_finite_overlap_norm_compression}
        \|P_iP_jv\|\le (1-\gamma)\frac{1}{m}\|P_iv\|.
    \end{equation}
    Then $H=\sum_iP_i$ satisfies $H^2\ge\gamma H$ as a quadratic form. The same
    conclusion holds from a separate compression coefficient $\eta$ whenever
    $\eta\le (1-\gamma)/m$. In particular, if $0\le \eta$ and $\eta m<1$, then
    \[
        H^2\ge (1-\eta m)H
    \]
    and $1-\eta m>0$.
\end{lemma}

\begin{proof}
    \leanok
    \uses{lem:finite_overlap_projection_rowsum_reduction,
        rem:projection_geometry_rowsum_identities}
    Cauchy--Schwarz and symmetric idempotence turn the norm-compression estimate
    into
    $\re\langle P_i v,P_jv\rangle
        \ge -(1-\gamma)m^{-1}\re\langle P_iv,v\rangle$.
    If a separate coefficient $\eta$ is used, the assumption
    $\eta\le(1-\gamma)/m$ first weakens that estimate to
    (\ref{eq:ph_finite_overlap_norm_compression}). The finite-overlap row-sum
    reduction then applies. For the strict form take $\gamma=1-\eta m$.
\end{proof}
